\hojaejercicio{Interpolación e integración numérica}

% Ejercicio 1
\ejercicio{
    Hallar el polinomio $P$ que interpola la función $f(x) = \cos(\pi x)$ en los puntos $x_0 = 0$, $x_1 = \dfrac{1}{2}$, $x_2 = 1$, $x_3 = \dfrac{3}{2}$.
}{
    Si tenemos la función $f(x) = \cos(\pi x)$ y los puntos $x_0 = 0$, $x_1 = \dfrac{1}{2}$, $x_2 = 1$, $x_3 = \dfrac{3}{2}$, entonces construimos la tabla de valores:
    \begin{center}
        \begin{tabular}{|c|cccc|}
            \hline
            $x$ & $0$ & $\dfrac{1}{2}$ & $1$ & $\dfrac{3}{2}$ \\
            \hline
            $y$ & $1$ & $0$ & $-1$ & $0$ \\
            \hline
        \end{tabular}
    \end{center}
    Ahora para construir el polinomio de interpolación, podemos usar la fórmula de Lagrange:
    \[
        P(x) = \sum_{i=0}^3 y_i L_i(x)
    \]
    Donde:
    \[
        L_i(x) = \prod_{j=0, j \neq i}^3 \frac{x - x_j}{x_i - x_j}
    \]
    Procedamos a calcular cada $L_i(x)$:
    \begin{itemize}
        \item Para $i=0$:
        \[
            L_0(x) = \frac{(x - \frac{1}{2})(x - 1)(x - \frac{3}{2})}{(0 - \frac{1}{2})(0 - 1)(0 - \frac{3}{2})} = \frac{(x - \frac{1}{2})(x - 1)(x - \frac{3}{2})}{\frac{3}{4}} = \frac{4}{3}(x - \frac{1}{2})(x - 1)(x - \frac{3}{2})
        \]
        \item Para $i=1$ no es necesario ya que:
        \[
            y_1 = 0 \implies y_1 L_1(x) = 0
        \]
        \item Para $i=2$:
        \[
            L_2(x) = \frac{(x - 0)(x - \frac{1}{2})(x - \frac{3}{2})}{(1 - 0)(1 - \frac{1}{2})(1 - \frac{3}{2})} = \frac{x(x - \frac{1}{2})(x - \frac{3}{2})}{-\frac{1}{4}} = -4x(x - \frac{1}{2})(x - \frac{3}{2})
        \]
        \item Para $i=3$ no es necesario ya que:
        \[
            y_3 = 0 \implies y_3 L_3(x) = 0
        \]
    \end{itemize}
    Dejando el polinomio:
    \[
        P(x) = 1 \cdot L_0(x) + (-1) \cdot L_2(x) = \frac{4}{3}(x - \frac{1}{2})(x - 1)(x - \frac{3}{2}) + 4x(x - \frac{1}{2})(x - \frac{3}{2})
    \]
}

% Ejercicio 2
\ejercicio{}{}

% Ejercicio 3
\ejercicio{}{}